% ============================================================================
% ANHANG K: BEC-Objekte – Die kompaktesten Strukturen im Universum
% ============================================================================

\chapter{BEC-Objekte – Die kompaktesten Strukturen im Universum}
\label{app:bec}

\section{Einleitung}
\label{app:bec_einleitung}

In der WDBT+ werden die kompaktesten Objekte des Universums nicht als 
Schwarze Löcher mit Singularitäten beschrieben, sondern als 
Bose-Einstein-Kondensate (BEC) aus Cooper-Paaren (siehe Kapitel~\ref{ch:bec_objekte}). 
Diese BEC-Objekte sind singularitätsfrei, ihre minimale Größe wird durch 
die fundamentale Längenskala \( l_0 \) begrenzt, und ihre maximale Masse 
folgt direkt aus der fraktalen Raumdimension \( D \).

Dieser Anhang liefert die vollständige, konsistente Herleitung aller 
relevanten Eigenschaften von BEC-Objekten.

\section{Die fundamentale Längenskala \( l_0 \)}
\label{app:bec_l0}

\subsection{Geometrischer Ursprung}
\label{sub:bec_l0_ursprung}

Aus der fraktalen Raumstruktur mit Dimension \( D \approx 2,704 \) und der 
Dodekaeder-Geometrie folgt (siehe Kapitel~\ref{ch:fraktale_dimension} und \ref{ch:bec_objekte}):

\[
l_0 = \left( \frac{V_{\text{Dodekaeder}}}{V_{\text{Einheitskugel}}} \right)^{1/3} \lambda_p \tag{K.1}
\]

Dabei ist \( \lambda_p = \hbar / (m_p c) \) die Compton-Wellenlänge des 
Protons (\( \lambda_p \approx 1,32 \times 10^{-15} \, \text{m} \)).

\subsection{Numerischer Wert}
\label{sub:bec_l0_numerisch}

Mit \( V_{\text{Dodekaeder}} / V_{\text{Einheitskugel}} \approx 2,5 \) folgt:

\[
l_0 \approx (2,5)^{1/3} \cdot 1,32 \times 10^{-15} \, \text{m} \approx 1,8 \times 10^{-15} \, \text{m} \tag{K.2}
\]

\subsection{Physikalische Bedeutung}
\label{sub:bec_l0_bedeutung}

Die Länge \( l_0 \) ist die kleinste physikalisch mögliche Distanz – 
die fundamentale Gitterkonstante des fraktalen Raumes.

\textbf{Wichtige Klarstellung:} \( l_0 \) ist \emph{nicht} die Compton-Wellenlänge 
des Neutrinos. Diese ist viel größer (siehe Anhang~\ref{app:neutrino_oszillation}):

\[
\lambda_\nu = \frac{\hbar}{m_\nu c} = l_0 \left( \frac{L_{Q,0}}{l_0} \right)^{(3-D)/2} \approx 2 \times 10^{-9} \, \text{m} \tag{K.3}
\]

mit \( L_{Q,0} = 10^{26} \, \text{m} \) (Hubble-Länge).

Die minimale Größe eines BEC-Objekts ist jedoch weiterhin \( R_{\min} = l_0 \), 
da BEC-Objekte aus Cooper-Paaren (gebundenen Fermionen) aufgebaut sind, 
nicht aus Neutrinos. Die unterschiedlichen Skalen sind physikalisch konsistent:
Das Neutrino ist nichtlokal und vermittelt das Quantenpotential über große 
Distanzen, während BEC-Objekte lokale materielle Kondensate sind.

\subsection{Die Neutrino-Compton-Wellenlänge im Vergleich}
\label{sub:bec_l0_compton_vergleich}

Die folgende Tabelle fasst die verschiedenen Skalen zusammen:

\[
\begin{array}{l|c|c}
\text{Größe} & \text{Symbol} & \text{Numerischer Wert} \\
\hline
\text{Fundamentale Gitterkonstante} & l_0 & 1,8 \times 10^{-15} \, \text{m} \\
\text{Compton-Wellenlänge des Neutrinos} & \lambda_\nu & 2 \times 10^{-9} \, \text{m} \\
\text{Compton-Wellenlänge des Protons} & \lambda_p & 1,32 \times 10^{-15} \, \text{m} \\
\text{Korrelationslänge des Q-Feldes} & L_{Q,0} & 10^{26} \, \text{m}
\end{array} \tag{K.4}
\]

Die Hierarchie \( l_0 \sim \lambda_p \ll \lambda_\nu \ll L_{Q,0} \) ist 
eine natürliche Konsequenz der fraktalen Geometrie und der nichtlokalen 
Natur des Q-Feldes.

\section{Paarbildung und Bose-Einstein-Kondensation}
\label{app:bec_paarbildung}

\subsection{Cooper-Paarbildung}
\label{sub:bec_cooper}

Bei hinreichend hohen Dichten bilden Fermionen Cooper-Paare und werden 
zu effektiven Bosonen. Die Paarbindungsenergie ist:

\[
E_{\text{pair}} = \frac{\hbar^2}{2m^* \xi^2} \tag{K.5}
\]

mit \( m^* \) der effektiven Masse und \( \xi \) der Kohärenzlänge.

\subsection{Bose-Einstein-Kondensation}
\label{sub:bec_kondensation}

Die kritische Temperatur für die Kondensation ist:

\[
T_c = \frac{2\pi \hbar^2}{k_B m^*} \left( \frac{n}{\zeta(3/2)} \right)^{2/3} \tag{K.6}
\]

Für typische Kernmaterie-Dichten (\( n \sim 10^{44} \, \text{m}^{-3} \)) 
liegt \( T_c \) im Bereich von \( 10^{12} \, \text{K} \).

\section{Zustandsgleichung eines BEC-Objekts im fraktalen Raum}
\label{app:bec_zustandsgleichung}

\subsection{Energiedichte}
\label{sub:bec_energie}

Die Gesamtenergie eines BEC-Objekts ist:

\[
E = E_{\text{kin}} + E_{\text{int}} + E_{\text{grav}} \tag{K.7}
\]

Im Thomas-Fermi-Limes dominiert die Wechselwirkungsenergie. 
Für ein BEC mit Kontaktwechselwirkung gilt:

\[
P(\rho) = \frac{g}{2} \rho^2 \tag{K.8}
\]

mit \( g = 4\pi\hbar^2 a_s / m_B \), \( a_s \) der Streulänge und 
\( m_B \approx 2m_p \).

\subsection{Fraktale Modifikation}
\label{sub:bec_fraktal}

Im fraktalen Raum mit Dimension \( D \) folgt aus der Gleichgewichtsbedingung 
zwischen Gravitation und Quantenpotential:

\[
P(\rho) = K \cdot \rho^{\gamma} \tag{K.9}
\]

mit

\[
\gamma = 1 + \frac{2}{D}, \quad 
K = \frac{\hbar^2}{2m_B} \left( \frac{3\pi^2}{g} \right)^{2/3} \cdot 
\frac{1}{m_B^{2/3}} \cdot f(D) \tag{K.10}
\]

Dabei ist \( f(D) \approx 1,2 \) für \( D \approx 2,704 \).

\section{Radius-Masse-Beziehung}
\label{app:bec_radius_masse}

\subsection{Gravitationsenergie im fraktalen Raum}
\label{sub:bec_grav_energie}

Für eine homogene Kugel vom Radius \( R \) und Gesamtmasse \( M \) 
skaliert die gravitative Selbstenergie wie:

\[
E_{\text{grav}} = -C_D \frac{GM^2}{R^{D-2}} \tag{K.11}
\]

Für \( D = 3 \) gilt \( C_3 = 3/5 \). Für \( D \approx 2,704 \) ist 
\( C_D \approx 0,4 \).

\subsection{Quantenpotential-Energie}
\label{sub:bec_q_energie}

Die Energie des Quantenpotentials skaliert wie:

\[
E_Q = A \frac{\hbar^2 M}{m_B^2 R^{2D/3}} \tag{K.12}
\]

mit \( A \approx 1 \).

\subsection{Minimierung der Gesamtenergie}
\label{sub:bec_minimierung}

Die Gesamtenergie ist:

\[
E(R) = -C_D \frac{GM^2}{R^{D-2}} + A \frac{\hbar^2 M}{m_B^2 R^{2D/3}} \tag{K.13}
\]

Ableitung nach \( R \) und Nullsetzen (\( dE/dR = 0 \)) liefert:

\[
C_D (D-2) \frac{GM^2}{R^{D-1}} = A \cdot \frac{2D}{3} \frac{\hbar^2 M}{m_B^2 R^{2D/3 + 1}} \tag{K.14}
\]

\subsection{Radius als Funktion der Masse}
\label{sub:bec_radius_formel}

Auflösen nach \( R \) ergibt:

\[
R^{2 - D/3} = \frac{2D A \hbar^2}{3(D-2) C_D G m_B^2} \cdot \frac{1}{M} \tag{K.15}
\]

Für \( D \approx 2,704 \) berechnet man:

\[
\frac{D}{3} = \frac{2,704}{3} = 0,90133
\]
\[
2 - \frac{D}{3} = 2 - 0,90133 = 1,09867
\]

Also:

\[
R^{1,09867} = \frac{2D A \hbar^2}{3(D-2) C_D G m_B^2} \cdot \frac{1}{M}
\]

Daraus folgt:

\[
R(M) = K_R \cdot M^{-1/1,09867} = K_R \cdot M^{-0,91} \tag{K.16}
\]

\subsection{Numerische Bestimmung von \( K_R \)}
\label{sub:bec_kr}

Die Konstante \( K_R \) folgt aus den Naturkonstanten:

\[
K_R = \left( \frac{2D A \hbar^2}{3(D-2) C_D G m_B^2} \right)^{1/1,09867} \tag{K.17}
\]

Einsetzen der Werte ergibt:

\[
K_R \approx 5,7 \times 10^{13} \, \text{m} \cdot \text{kg}^{0,91} \tag{K.18}
\]

\section{Maximale Masse eines BEC-Objekts}
\label{app:bec_max_masse}

\subsection{Untere Grenze: \( R_{\min} = l_0 \)}
\label{sub:bec_min_radius}

Die minimale physikalisch mögliche Größe ist \( l_0 \). 
Setzt man \( R = l_0 \) in (K.16) ein:

\[
M_{\text{BEC}} = \left( \frac{K_R}{l_0} \right)^{1/0,91} \tag{K.19}
\]

Mit \( K_R \approx 5,7 \times 10^{13} \, \text{m} \cdot \text{kg}^{0,91} \) 
und \( l_0 \approx 1,8 \times 10^{-15} \, \text{m} \):

\[
\frac{K_R}{l_0} \approx 3,2 \times 10^{28} \, \text{kg}^{0,91}
\]

Die Umrechnung in Sonnenmassen (\( 1 M_{\odot} \approx 2 \times 10^{30} \, \text{kg} \)) 
und die Potenz \( 1/0,91 \approx 1,099 \) liefert:

\[
M_{\text{BEC}} \approx 3 \times 10^6 M_{\odot} \tag{K.20}
\]

\subsection{Physikalische Bedeutung}
\label{sub:bec_max_bedeutung}

Dies ist die maximale Masse eines einzelnen BEC-Objekts. 
Oberhalb dieser Grenze wäre \( R < l_0 \), was physikalisch unmöglich ist. 

\textbf{Wichtige Klarstellung:} Die maximale Masse \( M_{\text{BEC}} \) 
hängt \emph{nicht} von der Neutrinomasse ab. Sie wird allein durch 
die fundamentale Längenskala \( l_0 \) und die Masse der Cooper-Paare 
\( m_B \approx 2m_p \) bestimmt. Die Korrektur der Neutrinomasse hat 
daher \emph{keine} Auswirkungen auf \( M_{\text{BEC}} \).

Beobachtete supermassereiche Objekte mit Massen oberhalb dieser Grenze 
(z. B. in sehr großen Galaxien bis \( 10^{10} M_{\odot} \)) sind daher 
keine einzelnen kompakten Objekte, sondern Systeme aus einem 
BEC-Kern und einer massiven Umgebung (siehe Kapitel~\ref{ch:kosmologie}).

\section{Strukturelle Analogie zum Atom}
\label{app:bec_atom_analogie}

\subsection{Atomare Skalierung}
\label{sub:bec_atom_tabelle}

\begin{center}
\begin{tabular}{lcc}
\hline
 & \textbf{Atom} & \textbf{BEC-Galaxie} \\
\hline
Kernradius & \( \sim 10^{-15} \, \text{m} \) & \( l_0 \sim 10^{-15} \, \text{m} \) \\
Hüllenradius & \( \sim 10^{-10} \, \text{m} \) & \( \sim 10^{21} \, \text{m} \) \\
Massenverhältnis & \( M_{\text{Kern}} \approx M_{\text{ges}} \) & \( M_{\text{Kern}} \ll M_{\text{Gal}} \) \\
\hline
\end{tabular}
\end{center}

\subsection{Fraktale Verbindung}
\label{sub:bec_fraktale_verbindung}

Die fraktale Raumdimension \( D \approx 2,704 \) verbindet beide Welten – 
sie bestimmt sowohl die Feinstrukturkonstante (Atom) als auch das 
Massenverhältnis \( M_{\text{Gal}} / M_{\text{Kern}} \approx 1000 \) 
(Galaxie, siehe Kapitel~\ref{ch:kosmologie}).

\section{Relativistische Korrekturen aus der Weber-Gravitation}
\label{app:bec_relativistik}

\subsection{Modifizierte Kreisbahngeschwindigkeit}
\label{sub:bec_kreisbahn}

Für eine kreisförmige Bahn liefert die Weber-Kraft mit \( \beta = 0,5 \) 
(siehe Kapitel~\ref{ch:dstt_weber_kraefte}):

\[
\frac{mv^2}{r} = \frac{GMm}{r^2} \left( 1 + \frac{v^2}{2c^2} \right) \tag{K.21}
\]

\subsection{ISCO-Bedingung}
\label{sub:bec_isco}

Auflösen nach \( v^2 \):

\[
v^2 = \frac{GM}{r} \left( 1 + \frac{GM}{2c^2 r} \right) \tag{K.22}
\]

Die Bedingung \( dv/dr = 0 \) ergibt:

\[
r_{\text{ISCO}} = \frac{12GM}{c^2} \tag{K.23}
\]

Für \( M = M_{\text{BEC}} \approx 3 \times 10^6 M_{\odot} \):

\[
r_{\text{ISCO}} \approx 5,3 \times 10^{10} \, \text{m} \approx 0,35 \, \text{AU}
\]

\section{Beobachtete supermassereiche Objekte}
\label{app:bec_beobachtung}

\subsection{Systeme aus Kern und Umgebung}
\label{sub:bec_systeme}

Beobachtete zentrale Objekte in Galaxien mit Massen bis \( 10^{10} M_{\odot} \) 
sind nach der WDBT+ keine einzelnen kompakten Objekte, sondern Systeme 
aus einem zentralen BEC-Kern (\( M \approx M_{\text{BEC}} \)) und einer 
massiven Umgebung (Akkretionsscheibe, Sternhaufen, Gas, Neutrinos).

\subsection{Massenverhältnis}
\label{sub:bec_massenverhaeltnis}

Aus der fraktalen Skalierung und der DSTT-Drift folgt (siehe Kapitel~\ref{ch:kosmologie} 
sowie Anhang~\ref{app:identitaetsbeziehung} für die Herleitung von \( \alpha_{\text{in}} \)):

\[
\frac{M_{\text{Gal}}}{M_{\text{Kern}}} \approx 1000 \tag{K.24}
\]

Mit \( M_{\text{Kern}} = M_{\text{BEC}} \approx 3 \times 10^6 M_{\odot} \) 
ergibt sich \( M_{\text{Gal}} \approx 3 \times 10^9 M_{\odot} \), 
was typisch für große Galaxien ist.

\section{Kollisionen von BEC-Objekten}
\label{app:bec_kollisionen}

BEC-Objekte haben keine Sonderstellung. Bei Kollisionen verhalten sie sich 
wie normale physikalische Objekte:

\begin{itemize}
    \item \textbf{Verschmelzung:} Wenn \( M_1 + M_2 \leq M_{\text{BEC}} \), 
          können sie zu einem neuen BEC-Objekt verschmelzen.
    \item \textbf{Abprallen:} Wenn die Relativgeschwindigkeit hoch genug ist, 
          prallen sie elastisch auseinander.
    \item \textbf{Zerrüttung:} Bei geringer Relativgeschwindigkeit und 
          \( M_1 + M_2 > M_{\text{BEC}} \) können die Kerne zerreißen; 
          die Masse wird in die Umgebung überführt.
    \item \textbf{Doppelkern-Systeme:} Die Kerne bleiben getrennt und 
          umkreisen sich in einer gemeinsamen Umgebung.
\end{itemize}

Gravitationswellen von BEC-Kollisionen unterscheiden sich geringfügig 
von denen Schwarzer Löcher: Sie haben keinen Ereignishorizont, eine 
Oberfläche und möglicherweise ein Nachglühen (siehe Kapitel~\ref{ch:bec_objekte}).

\section{Verbindung zur Neutrinomasse}
\label{app:bec_neutrino_verbindung}

Die fundamentale Längenskala \( l_0 \) bestimmt nicht nur die minimale 
Größe von BEC-Objekten, sondern auch – über die Korrelationslänge 
\( L_{Q,0} = 10^{26} \, \text{m} \) – die Masse des Neutrinos 
(siehe Anhang~\ref{app:neutrino_oszillation}):

\[
m_\nu = \frac{\hbar}{l_0 c} \left( \frac{l_0}{L_{Q,0}} \right)^{(3-D)/2} \approx 0,1 \, \text{eV}/c^2
\]

Die Konsistenz beider Ableitungen – der BEC-Maximalmasse und der 
Neutrinomasse – aus derselben fundamentalen Längenskala \( l_0 \) 
und derselben fraktalen Dimension \( D \) ist ein starkes Indiz 
für die Richtigkeit der Theorie.

Die folgende Tabelle fasst die Hierarchie der Skalen zusammen:

\[
\begin{array}{l|c|c}
\text{Phänomen} & \text{Relevante Skala} & \text{Größenordnung} \\
\hline
\text{BEC-Objekt (minimale Größe)} & l_0 & 10^{-15} \, \text{m} \\
\text{Neutrino (Compton-Wellenlänge)} & \lambda_\nu & 10^{-9} \, \text{m} \\
\text{Q-Feld (Korrelationslänge)} & L_{Q,0} & 10^{26} \, \text{m} \\
\end{array}
\]

Die unterschiedlichen Skalen sind physikalisch konsistent: 
Das Neutrino ist nichtlokal und vermittelt das Quantenpotential 
über große Distanzen, während BEC-Objekte lokale materielle 
Kondensate aus Cooper-Paaren sind.

\section{Zusammenfassung}
\label{app:bec_zusammenfassung}

\begin{itemize}
    \item Minimale Größe: \( l_0 \approx 1,8 \times 10^{-15} \, \text{m} \) 
          (fundamentale Gitterkonstante des fraktalen Raumes)
    \item Radius-Masse-Beziehung: \( R \propto M^{-0,91} \)
    \item Maximale Masse: \( M_{\text{BEC}} \approx 3 \times 10^6 M_{\odot} \)
          (unabhängig von der Neutrinomasse)
    \item Beobachtete supermassereiche Objekte sind Systeme aus Kern und Umgebung
    \item Massenverhältnis Galaxie zu Kern: etwa 1000
    \item BEC-Objekte verhalten sich bei Kollisionen wie normale Objekte
    \item Die Compton-Wellenlänge des Neutrinos ist \( \lambda_\nu \approx 2 \times 10^{-9} \, \text{m} \) 
          – viel größer als \( l_0 \)
    \item Die gleiche fundamentale Längenskala \( l_0 \) bestimmt sowohl 
          die minimale BEC-Größe als auch – über \( L_{Q,0} \) – die Neutrinomasse
\end{itemize}

\section{Ausblick}
\label{app:bec_ausblick}

Die Vorhersage einer maximalen Masse \( M_{\text{BEC}} \approx 3 \times 10^6 M_{\odot} \) 
ist testbar (siehe Kapitel~\ref{ch:testbare_vorhersagen}). Zukünftige Beobachtungen mit dem 
Event Horizon Telescope (EHT) könnten die Unterscheidung zu Schwarzen Löchern 
ermöglichen:

\begin{itemize}
    \item Bei Galaxien mit zentralen Massen \( > 3 \times 10^6 M_{\odot} \) 
          sollte kein einfacher Schatten eines einzelnen kompakten Objekts 
          sichtbar sein, sondern eine komplexe Struktur.
    \item Die Feinstruktur des Schattens, Polarisationsmuster und 
          zeitliche Variabilität (Flares) unterscheiden sich von 
          ART-Vorhersagen (siehe Kapitel~\ref{ch:kosmologie}).
\end{itemize}

Die konsistente Erklärung der BEC-Maximalmasse, der Neutrinomasse und 
der Hubble-Spannung aus derselben fraktalen Geometrie (\( D \approx 2,704 \)) 
und derselben Korrelationslänge (\( L_{Q,0} = 10^{26} \, \text{m} \)) 
macht die WDBT+ zu einer geschlossenen, parameterfreien Theorie 
der fundamentalen Physik.